% Exam Template for UMTYMP and Math Department courses
%
% Using Philip Hirschhorn's exam.cls: http://www-math.mit.edu/~psh/#ExamCls
%
% run pdflatex on a finished exam at least three times to do the grading table on front page.
%
%%%%%%%%%%%%%%%%%%%%%%%%%%%%%%%%%%%%%%%%%%%%%%%%%%%%%%%%%%%%%%%%%%%%%%%%%%%%%%%%%%%%%%%%%%%%%%

% These lines can probably stay unchanged, although you can remove the last
% two packages if you're not making pictures with tikz.
\documentclass[11pt]{exam}
\RequirePackage{amssymb, amsfonts, amsmath, latexsym, verbatim, xspace, setspace}
\RequirePackage{tikz, pgflibraryplotmarks}

% By default LaTeX uses large margins.  This doesn't work well on exams; problems
% end up in the "middle" of the page, reducing the amount of space for students
% to work on them.
\usepackage[margin=1in]{geometry}
\usepackage{enumerate}
\usepackage{amsthm}

\theoremstyle{definition}
\newtheorem{soln}{Solution}

% Here's where you edit the Class, Exam, Date, etc.
\newcommand{\class}{Math 150A Section 8 and 16}
\newcommand{\term}{Spring 2025}
\newcommand{\examnum}{Exam II}
\newcommand{\examdate}{March 13, 2025}
\newcommand{\timelimit}{1 Hour 50 Minutes}

% For an exam, single spacing is most appropriate
\singlespacing
% \onehalfspacing
% \doublespacing

% For an exam, we generally want to turn off paragraph indentation
\parindent 0ex

\begin{document} 

% These commands set up the running header on the top of the exam pages
\pagestyle{head}
\firstpageheader{}{}{}
\runningheader{\class}{\examnum\ - Page \thepage\ of \numpages}{\examdate}
\runningheadrule

\begin{flushright}
\begin{tabular}{p{2.8in} r l}
\textbf{\class} & \textbf{Name (Print):} & \makebox[2in]{\hrulefill}\\
\textbf{\term} &&\\
\textbf{\examnum} & \textbf{Student ID:}&\makebox[2in]{\hrulefill}\\
\textbf{\examdate} &&\\
\textbf{Time Limit: \timelimit} % & Teaching Assistant & \makebox[2in]{\hrulefill}
\end{tabular}\\
\end{flushright}
\rule[1ex]{\textwidth}{.1pt}


This exam contains \numpages\ pages (including this cover page) and
\numquestions\ problems.  Check to see if any pages are missing.  Enter
all requested information on the top of this page, and put your initials
on the top of every page, in case the pages become separated.\\

You may \textit{not} use your books or notes on this exam.  However, you may use a single, handwritten, one-sided notesheet and a \textit{basic} calculator.\\

You are required to show your work on each problem on this exam.  The following rules apply:\\

\begin{minipage}[t]{3.7in}
\vspace{0pt}
\begin{itemize}

%\item \textbf{If you use a ``fundamental theorem'' you must indicate this} and explain
%why the theorem may be applied.

\item \textbf{Organize your work}, in a reasonably neat and coherent way, in
the space provided. Work scattered all over the page without a clear ordering will 
receive very little credit.  

\item \textbf{Mysterious or unsupported answers will not receive full
credit}.  A correct answer, unsupported by calculations, explanation,
or algebraic work will receive no credit; an incorrect answer supported
by substantially correct calculations and explanations might still receive
partial credit.  This especially applies to limit calculations.

\item If you need more space, use the back of the pages; clearly indicate when you have done this.

\item \textbf{Box Your Answer} where appropriate, in order to clearly indicate what you consider the answer to the question to be.
\end{itemize}

Do not write in the table to the right.
\end{minipage}
\hfill
\begin{minipage}[t]{2.3in}
\vspace{0pt}
%\cellwidth{3em}
\gradetablestretch{2}
\vqword{Problem}
\addpoints % required here by exam.cls, even though questions haven't started yet.	
\gradetable[v]%[pages]  % Use [pages] to have grading table by page instead of question

\end{minipage}
\newpage % End of cover page

%%%%%%%%%%%%%%%%%%%%%%%%%%%%%%%%%%%%%%%%%%%%%%%%%%%%%%%%%%%%%%%%%%%%%%%%%%%%%%%%%%%%%
%
% See http://www-math.mit.edu/~psh/#ExamCls for full documentation, but the questions
% below give an idea of how to write questions [with parts] and have the points
% tracked automatically on the cover page.
%
%
%%%%%%%%%%%%%%%%%%%%%%%%%%%%%%%%%%%%%%%%%%%%%%%%%%%%%%%%%%%%%%%%%%%%%%%%%%%%%%%%%%%%%

\begin{questions}

\addpoints

\question[10]\mbox{}
\textbf{TRUE or FALSE!}  Write 
\begin{enumerate}[(a)]
\item every inflection point of $f(x)$ is also a critical point of $f(x)$
\vspace{1.2in}
\item the function $f(x) = x^{1/3}$ has no critical points
\vspace{1.2in}
\item if $f(x)$ has three different $x$-intercepts, then $f'(x)$ must have at least two $x$-intercepts
\vspace{1.2in}
\item every function has an absolute maximum
\vspace{1.2in}
\item if $f''(3) = -2$, then $f(x)$ has a local maximum at $x=3$
\vspace{1.2in}
\end{enumerate}

\newpage
\question[10]\mbox{}

We want to construct a cylindrical can with a bottom but no top that will have a volume of 30 cm3. Determine the dimensions of the can that will minimize the amount of material needed to construct the can.

Hint: the volume of a cylinder is $\pi r^2h$, the circumference of a circle is $2\pi r$ and the area of a circle is $\pi r^2$.

\newpage
\question[10]

Consider the function
$$g(t) = t^5-5t^4+8$$.

\begin{enumerate}[(a)]
\item Identify the critical points of the function
\item Determine the intervals on which the function increases and decreases.
\item Classify the critical points as relative maximums, relative minimums or neither.
\item Determine the intervals on which the function is concave up and concave down.
\item Determine the inflection points of the function.
\item Use the information from steps (a) – (e) to sketch the graph of the function.
\end{enumerate}

\newpage
\question[10]

Use L'Hospital's Rule to calculate the following limits:

\begin{enumerate}[(a)]
\item
$$\lim_{x\rightarrow 0} \frac{1}{x^2}-\frac{\cot(x)}{x}.$$
\item
$$\lim_{x\rightarrow\infty}x^{1/\ln(x)}.$$
\end{enumerate}

\newpage
\question[10]\mbox{}
Consider the equation
$$x^9 + 7x^5 + x^5 + 2x^3 + 3x + 4 = 0.$$
\begin{enumerate}[(a)]
\item  State the Mean Value Theorem
\vspace{2in}
\item  Use the Intermediate Value Theorem to show that the above equation has at least one solution
\vspace{3in}
\item  Use the Mean Value Theorem to explain why the equation has no more than one solution
\end{enumerate}

\newpage
\question[10]\mbox{}
For each of the following, provide an example or explain why it does not exist.
\begin{enumerate}[(a)]
\item  a function with no critical points
\vspace{1.2in}
\item  a function with inflection points at $-1$ and $1$
\vspace{1.2in}
\item  a nonconstant function $f(x)$ whose derivative is always zero
\vspace{1.2in}
\item  a function with a global maximum at $x=1$
\vspace{1.2in}
\item  a critical point of the function $3-|x-3|$
\vspace{1.2in}
\end{enumerate}

\newpage
\question[10]\mbox{}
Find the critical points of each of the following functions.
Classify them as local maxima, local minima, or neither.
\begin{enumerate}[(a)]
\item  $2x^3+9x^2+3x+4$
\vspace{4in}
\item  $9x^{5/3}-25x^{3/5}$
\end{enumerate}

\end{questions}
\end{document}
